\documentclass[11pt]{article}

\usepackage[T1]{fontenc}
\usepackage[utf8]{inputenc}
\usepackage{lmodern}
\usepackage[margin=1in]{geometry}
\usepackage{microtype}
\usepackage{hyperref}
\usepackage{enumitem}
\usepackage{titlesec}

\hypersetup{
  colorlinks=true,
  linkcolor=black,
  citecolor=blue,
  urlcolor=blue
}

\setlength{\parindent}{1.5em}
\setlength{\parskip}{0.25em}
\setlist[itemize]{leftmargin=1.5em,topsep=0.25em,itemsep=0.12em}
\titleformat{\section}{\normalfont\bfseries\large}{\thesection.}{0.6em}{}
\titleformat{\subsection}{\normalfont\bfseries\normalsize}{\thesubsection.}{0.6em}{}

\title{\textbf{OpenScientists: Public Infrastructure for AI-Assisted Discovery}\\[0.4em]
\large Proposal Summary}
\author{}
\date{April 2026}

\begin{document}
\maketitle

\begin{abstract}
OpenScientists is a public-interest initiative to build open infrastructure for AI-assisted scientific discovery. The core premise is that, if AI systems become deeply integrated into scientific workflows, the governing infrastructure should not be controlled exclusively by a small number of private platforms. We propose a coordinated program spanning tacit-knowledge capture, agent runtime systems, benchmark design, model development, and metascience.
\end{abstract}

\section{Motivation}
Scientific progress is a leverage point on nearly every other domain of human advancement. Improvements in the process of discovery propagate into energy, medicine, materials, computation, and climate. As AI systems become capable of assisting with literature synthesis, hypothesis generation, experiment design, coding, and analysis, they may substantially alter the rate and character of scientific production.

At present, much of this infrastructure is being developed inside a small number of frontier laboratories and commercial platforms. If the default tooling, evaluation standards, and workflow patterns become locked into proprietary systems before credible public alternatives exist, the future direction of AI-assisted science may be shaped primarily by commercial incentives. OpenScientists is motivated by the view that scientific infrastructure should remain inspectable, contestable, and aligned with the norms of reproducibility and open inquiry.

\section{Goal, Output, and Milestone}
The goal of OpenScientists is to build a credible open ecosystem for AI-assisted science: one that can support real researchers, produce measurable scientific value, and remain legible to public institutions rather than only to private frontier labs.

The intended outputs and milestones include:
\begin{itemize}
    \item \textbf{Tacit knowledge library:} a structured, versioned repository of scientific heuristics, failure modes, debugging instincts, and practical judgment that rarely appears in published papers.
    \item \textbf{Agent runtime system:} an execution environment for AI-assisted research, including literature search, tool use, code execution, persistent memory, and citation tracing.
    \item \textbf{Scientific ability benchmark:} evaluation based on genuine research tasks such as hypothesis generation, confounder detection, ablation design, methodological critique, and uncertainty calibration.
    \item \textbf{Model development:} open model improvement for scientific reasoning, emphasizing calibrated uncertainty, disciplined tool use, and long-horizon planning rather than fluent text alone.
    \item \textbf{Metascience and governance:} empirical study of how these systems affect scientific production, reviewer burden, citation networks, publication quality, and institutional norms.
    \item \textbf{10{,}000 ICLR papers:} a large, open, researcher-facing corpus and workflow testbed for literature-grounded agents, benchmark tasks, retrieval, citation tracing, and metascientific analysis.
\end{itemize}

Taken together, these outputs are meant to form a single public stack rather than isolated tools. Tacit knowledge improves agent behavior, runtime systems make models usable in real research practice, benchmarks create meaningful feedback loops, and metascience measures whether the system improves science rather than merely accelerating text production.

\section{Resources and Partnerships}
The initiative requires coordinated support across several categories:

\begin{itemize}
    \item \textbf{Research contributors:} PhD students and researchers willing to contribute tacit knowledge, domain heuristics, and failure cases.
    \item \textbf{Scientific users and testers:} researchers who deploy the tools in real workflows and provide grounded feedback.
    \item \textbf{Senior reviewers:} domain experts who can validate quality and distinguish substantive insight from superficial output.
    \item \textbf{Financial sponsors:} philanthropic or mission-aligned funders supporting engineering, fellowships, benchmark prizes, and long-term maintenance.
    \item \textbf{API and token sponsors:} model-provider credits for benchmarking and multi-model evaluation.
    \item \textbf{GPU and compute sponsors:} hardware and cloud support for training and large-scale evaluation.
    \item \textbf{Institutional partners:} universities, publishers, and scientific societies supporting pilot deployments, benchmark design, and governance experimentation.
\end{itemize}

\section{Why This Matters Now}
The timing is important. AI-assisted research systems have already begun to enter real workflows. Once institutions standardize on closed platforms, the costs of building a credible open alternative increase substantially. Standards, habits, and interfaces harden quickly.

For that reason, OpenScientists should be understood not as a general statement in favor of openness, but as a time-sensitive infrastructure effort. If open scientific AI is to remain viable, its foundations must be built while the design space is still fluid.

\end{document}
